\documentclass[10pt,a4paper]{amsart}
\usepackage[margin=19mm]{geometry}
\usepackage{amsmath,amssymb,mathtools}
\usepackage[scr=rsfs,cal=euler]{mathalfa}
\usepackage{xfrac}
\usepackage[unicode,hidelinks,bookmarks=false]{hyperref}
\newcommand{\ie}{i.e.\ }
\newcommand{\cf}{cf.\ }
\newcommand{\resp}{respectively\ }
\newcommand{\Subsection}{\subsection}
\providecommand{\nobreakdash}{\nobreak\hbox{-}}
\setlength{\emergencystretch}{2em}
\multlinegap0pt
\usepackage{amssymb}
\usepackage{tikz}
\usepackage{mathtools}
\newtheorem{cor}{Corollary}
\newcommand{\C}{\mathcal{C}}
\newcommand{\F}{\mathcal{F}}
\newcommand{\conv}{\operatorname{conv}}
\title{Strong invariants and Tverberg numbers in convexity spaces}
\author{Minho Cho, Andreas F. Holmsen, Attila Jung, Hong Liu}
\date{Source: arXiv:2607.29403v1 (CC BY 4.0)}
\begin{document}
\maketitle

\noindent\emph{Source and licence.} Strong invariants and Tverberg numbers in convexity spaces. Version arXiv:2607.29403v1 (CC BY 4.0), distributed by arXiv under the Creative Commons Attribution 4.0 International licence, http://creativecommons.org/licenses/by/4.0/, as recorded in the arXiv licence metadata for this exact version and shown on the abstract page https://arxiv.org/abs/2607.29403v1. Full licence notice: \texttt{licenses/arxiv-2607.29403-CC-BY-4.0.txt}.\par\smallskip

\noindent\textit{Editorial scope.} Counted unit: the article's cor cor:separable-pq (source0.tex lines 1483--1498). The article's complete original proof of it is delivered with it (source0.tex lines 1500--1558, one \texttt{proof} environment of 59 lines). No mathematics is written, repaired or re-derived: the statement and the proof are the article's own lines, in the article's own notation, and no retained line is substituted or re-flowed.  Retained uncounted as the source-local context the proof needs: lines 1438--1446.  Nothing was omitted from any retained span, so the package carries no omission record.  No retained line contains a citation, so the wrapper carries no bibliography and no bibliography entry.  No retained line contains a figure environment, an image-inclusion command or drawing code, so no image is delivered and the package declares no asset dependency.  Editorial formatting only, in addition: an AMS \texttt{amsart} preamble that restates the article's own declarations and macros used by the retained lines, namely the environments \texttt{cor} on separate counters, together with the standard packages those lines need.  The retained environments print in this document as Corollary 1 in order of appearance, whereas the article prints the same items with the numbering of its own full document.  The complete native source of the article as distributed in the arXiv e-print for this exact version is bundled unchanged as \texttt{sources/206480/source0.tex}.
\begin{cor}[Weak $\varepsilon$-nets]
\label{cor:separable-weak-epsilon-net}
For every finite point set or labelled point multiset $P$ and every
$0<\varepsilon\le1$, there is a weak $\varepsilon$-net
$N\subseteq X$ satisfying
\[
|N|\le 2\left(\frac e\varepsilon\right)^a.
\]
\end{cor}

\begin{cor}[A quantitative $(p,q)$-theorem]
\label{cor:separable-pq}
Assume that $\C$ has fractional Helly number at most $q_0$ with
fractional Helly function $\beta$.  Let $p\ge q\ge q_0$, and let
$\F\subseteq\C$ be a finite family of nonempty convex sets satisfying
the $(p,q)$-property.  Put
\[
p'=(p-1)(q-1)+1,
\qquad
\alpha=\frac{\binom q{q_0}}{\binom{p'}{q_0}},
\qquad
\gamma=\beta(\alpha).
\]
Then $\F$ has a transversal of size at most
$2\left(\frac e\gamma\right)^a.$
\end{cor}

\begin{proof}
We first prove a weighted consequence of fractional Helly.  Let
$\F'$ be an arbitrary finite labelled multiset whose members are
taken from $\F$.  Repeating all labels equally if necessary, assume
that $n=|\F'|\ge p'$.

Every $p'$ labelled members of $\F'$ contain $q$ members with
nonempty intersection.  Indeed, either they contain at least $p$
distinct members of $\F$, in which case the $(p,q)$-property applies,
or some member occurs at least
$\left\lceil\frac{p'}{p-1}\right\rceil=q$
times.  Let $M_{q_0}$ be the number of intersecting labelled
$q_0$-subfamilies of $\F'$.  Counting pairs $(I,J)$, where $J$ is a
$p'$-subfamily and $I\subseteq J$ is an intersecting
$q_0$-subfamily, gives
$M_{q_0}\binom{n-q_0}{p'-q_0}
\ge
\binom n{p'}\binom q{q_0}.$
Since
$\binom n{p'}\binom{p'}{q_0}
=
\binom n{q_0}\binom{n-q_0}{p'-q_0},$
we obtain
$\frac{M_{q_0}}{\binom n{q_0}}
\ge
\frac{\binom q{q_0}}{\binom{p'}{q_0}}
=
\alpha.$
The fractional Helly property therefore yields a point contained in
at least $\gamma n$ members of $\F'$, counted with multiplicity.

It follows, by approximation, that for every probability distribution
on the members of $\F$, some point is contained in sets of total
weight at least $\gamma$.  There are only finitely many incidence
patterns of points on the finite family $\F$.  Applying the minimax
theorem to the resulting finite zero--one matrix gives a finitely
supported probability distribution $\mu$ on $X$ such that
\[
\mu(F)\ge\gamma
\qquad\text{for every }F\in\F.
\]
Equivalently, this is the usual finite-dimensional linear-programming
duality argument.  Since the corresponding finite game has a rational
optimal strategy, we may clear denominators and obtain a finite
labelled multiset $Y$ of points satisfying
\begin{equation}\label{eq:fractional-piercing-measure}
|F\cap Y|\ge\gamma|Y|
\qquad\text{for every }F\in\F.
\end{equation}

Apply Corollary~\ref{cor:separable-weak-epsilon-net} to $Y$ with
$\varepsilon=\gamma$.  We obtain a set $N\subseteq X$ of size at most
$2(e/\gamma)^a$ which meets the convex hull of every submultiset of
$Y$ of size at least $\gamma|Y|$.  For every $F\in\F$,
\eqref{eq:fractional-piercing-measure} and the convexity of $F$ give
$\conv(F\cap Y)\subseteq F.$
Hence $N\cap F\ne\varnothing$ for every $F\in\F$, so $N$ is the
required transversal.
\end{proof}

\end{document}
